% Part: quantified-modal-logic
% Chapter: semantics
% Section: satisfaction-free

\documentclass[../../../include/open-logic-section]{subfiles}

\begin{document}

\olfileid{qml}{sem}{saf}

\olsection{Satisfaction in Variable-Domain Models}

\begin{explain}
The definition of !!^{variable} assignments is as in fixed-domain models except that we use 
the super-domain instead of the single fixed domain. 
\end{explain}

\begin{defn}[Variable Assignment]
A \emph{variable assignment}~$s$ for a variable-domain model~$\mModel{M}$ 
is a function which maps each !!{variable} to !!a{element} of
the super-domain~$\mathcal{S}$,
i.e., $s\colon \Var \to \mathcal{S}$.
\end{defn}

\begin{explain}
The definitions of terms !!{value}s and of 
assignment variants are unchanged.
\end{explain}

\begin{defn}[!!^{value} of Terms]
If $t$ is a term of the language~$\Lang L$, $\mModel M$ is a
variable-domain model for~$\Lang L$, $w$ a world in $mWorlds M$, and $s$
is !!a{variable} assignment for~$\mModel M$, the
\emph{!!{value}}~$\mValue{t}{M}{w}[s]$ of $t$ at world $w$ in the is
defined as follows:
\begin{enumerate}
\item \indcase{t}{c}{$\mValue{\indfrm}{M}{w}[s] = \mAssign{\indcomplex}{M}$.}
\item \indcase{t}{x}{$\mValue{\indfrm}{M}{w}[s] = s(\indcomplex)$.}
\item \indcase{t}{\Atom{f}{t_1, \ldots, t_n}}{
\[
\mValue{\indfrm}{M}{w}[s] = \mAssign{f}{M}[w](\mValue{t_1}{M}{w}[s], \ldots,
\mValue{t_n}{M}{w}[s]).
\]}
\end{enumerate}
\end{defn}

\begin{defn}[$x$-Variant]
If $s$ is !!a{variable} assignment for a variable-domain
model~$\mModel{M}$, then any !!{variable} assignment~$s'$ for~$\mModel
M$ which differs from~$s$ at most in what it assigns to~$x$ is called
an \emph{$x$-variant} of~$s$. If $s'$ is an $x$-variant of~$s$ we
write $\varAssign{s'}{s}{x}$.
\end{defn}

\begin{defn}
  If $s$ is !!a{variable} assignment for a variable-domain model~$\mModel{M}$, 
  and $m \in \mDomain{M}$, then the assignment~$\Subst{s}{m}{x}$ is the
  variable assignment defined by
  \[\Subst{s}{m}{x}(y) = \begin{cases}
    m & \text{if } y \ident x\\
    s(y) & \text{otherwise}.
  \end{cases}\]
\end{defn}

\subsection{Satisfaction}

\begin{explain}

Satisfaction is defined as for fixed-domain models, with two differences:
\begin{enumerate}
\item For quantifiers: make their range world-relative,
\item For identity, in the negative semantics: restrict its application
to each world's domain.
\end{enumerate}
Let us look at these differences before stating the full 
definitions of satisfactions. 

\paragraph{Quantifiers} With variable-domain models we restrict the range of 
quantifiers at each world to that world's domain. That includes
the existence predicate:
\begin{itemize}
\item $\qmSat{M}{\Atom{\lfrexists}{t}}[w][s]$ iff 
  $\mValue{t}{M}{w}[s]\in \mDomain{M}(w)$.

\tagitem{prvAll}{%
  $\qmSat{M}{\lforall[x][!B]}[w][s]$ iff for
    every !!{element}~$m \in \mDomain{M}[w]$,
    $\qmSat{M}{!B}[w][\Subst{s}{m}{x}]$.}{}

\tagitem{prvEx}{%
  $\qmSat{M}{\lexists[x][!B]}[w][s]$ iff for at
  least one !!{element}~$m \in \mDomain{M}[w]$,
  $\qmSat{M}{!B}[w][\Subst{s}{m}{x}]$.}{}
\end{itemize}
Note that for a quantified sentence to hold at a world 
$w$ relative to an assignment $s$, we only variants
of $s$ involving objects \emph{of that world's domain}, 
$\mDomain{M}[w]$---not objects of the model's super-domain
 $\mathcal{S}$.

\paragraph{Identity in Negative Free semantics}. For the \emph{positive}
variant of quantified modal logic, the clause for identity
is the same. For the \emph{negative} variant, we restrict it 
as follows:

\begin{itemize}

\item classical, positive free: $\qmSat{M}{\eq[t_1][t_2]}[w][s]$ iff
  $\mValue{t_1}{M}{w}[s] = \mValue{t_2}{M}{w}[s]$.

\item negative free: $\qmSat{M}{\eq[t_1][t_2]}[w][s]$ iff
  $\mValue{t_1}{M}{w}[s]\in \mDomain{M}[w]$ and 
  $\mValue{t_1}{M}{w}[s] = \mValue{t_2}{M}{w}[s]$.

\end{itemize}

\end{explain}


% TO DO:   
\begin{editorial}
  Hard-coding $\models$ rather than using `Sat' here. 
  Need to provide a way to index the satisfaction to a given logic?
\end{editorial}

Satisfaction has a different identity clause for the 
positive and for the negative variants of free quantified 
modal logic. When we need to disambiguate we write
 $\models_{\Log{PFQML}}$ for the positive semantics and
 $\models_{\Log{NFQML}}$. 

\begin{defn}[Satisfaction in variable-domain models]
\ollabel{defn:satisfaction}
Satisfaction of !!a{formula}~$!A$ a fixed-domain model~$\mModel{M}$
at a world $w$ in $\mWorlds{M}$ relative to !!a{variable} assignment~$s$, in symbols:
$\qmSat{M}{!A}[w][s]$, is defined recursively as follows. We write
$\qmSat/{M}{!A}[w][s]$ to mean ``not $\qmSat{M}{!A}[w][s]$.'' 
When we need to specify the positive or negative variant of 
the definition, we use  When we need to disambiguate we write
 $\models_{\Log{PFQML}}$ for the positive semantics and
 $\models_{\Log{NFQML}}$. 
\begin{enumerate}
\tagitem{prvFalse}{%
  \indcase{!A}{\lfalse}{$\qmSat/{M}{\indfrm}[w][s]$.}}{}

\tagitem{prvTrue}{%
  \indcase{!A}{\ltrue}{$\qmSat{M}{\indfrm}[w][s]$.}}{}

\item \indcase{!A}{\Atom{R}{t_1, \dots, t_n}}{$\qmSat{M}{\indfrm}[w][s]$
  iff $\langle \mValue{t_1}{M}{w}[s], \dots, \mValue{t_n}{M}{w}[s] \rangle \in
  \mAssign{R}{M}[w]$.}

\item \indcase{!A}{\Atom{\lfrexists}{t}}{$\qmSat{M}{\indfrm}[w][s]$ iff 
  $\mValue{t}{M}{w}[s]\in \mDomain{M}(w)$.}

\item \indcase{!A}{\eq[t_1][t_2]}{$\qmSat{M}{\indfrm}[w][s]$ iff
  \begin{itemize}
    \item positive variant: $\mValue{t_1}{M}{w}[s] = \mValue{t_2}{M}{w}[s]$,
    \item negative variant:  $\mValue{t_1}{M}{w}[s]\in \mDomain{M}[w]$ and 
  $\mValue{t_1}{M}{w}[s] = \mValue{t_2}{M}{w}[s]$.
  \end{itemize}
 }

\tagitem{prvNot}{%
  \indcase{!A}{\lnot !B}{$\qmSat{M}{\indfrm}[w][s]$ iff
    $\qmSat/{M}{!B}[w][s]$.}}{}

\tagitem{prvAnd}{%
  \indcase{!A}{(!B \land !C)}{$\qmSat{M}{\indfrm}[w][s]$ iff 
  $\qmSat{M}{!B}[w][s]$ and $\qmSat{M}{!C}[w][s]$.}}{}

\tagitem{prvOr}{%
  \indcase{!A}{(!B \lor !C)}{$\qmSat{M}{\indfrm}[w][s]$ iff
    $\qmSat{M}{!B}[w][s]$ or $\qmSat{M}{!C}[w][s]$ (or both).}}{}

\tagitem{prvIf}{%
  \indcase{!A}{(!B \lif !C)}{$\qmSat{M}{\indfrm}[w][s]$ iff 
    $\qmSat/{M}{!B}[w][s]$ or $\qmSat{M}{!C}[w][s]$ (or both).}}{}

\tagitem{prvIff}{%
  \indcase{!A}{(!B \liff !C)}{$\qmSat{M}{\indfrm}[w][s]$ iff either
    both $\qmSat{M}{!B}[w][s]$ and $\qmSat{M}{!C}[w][s]$, or neither
    $\qmSat{M}{!B}[w][s]$ nor $\qmSat{M}{!C}[w][s]$.}}{}

\tagitem{prvAll}{%
  \indcase{!A}{\lforall[x][!B]}{$\qmSat{M}{\indfrm}[w][s]$ iff for
    every !!{element}~$m \in \mDomain{M}(w)$,
    $\qmSat{M}{!B}[w][\Subst{s}{m}{x}]$.}}{}

\tagitem{prvEx}{%
  \indcase{!A}{\lexists[x][!B]}{$\qmSat{M}{\indfrm}[w][s]$ iff for at
  least one !!{element}~$m \in \mDomain{M}(w)$,
  $\qmSat{M}{!B}[w][\Subst{s}{m}{x}]$.}}{}

\tagitem{prvBox}{%
  \indcase{!A}{\Box !B}{$\qmSat{M}{\indfrm}[w][s]$ iff
    $\qmSat{M}{!B}[w'][s]$ for all $w' \in \mWorlds{M}$ with
    $\mAccess{M}ww'$.}}{}

\tagitem{prvDiamond}{%
    \indcase{!A}{\Diamond !B}{$\qmSat{M}{\indfrm}[w][s]$ iff
      $\qmSat{M}{!B}[w'][s]$ for at least one $w' \in \mWorlds{M}$
      with $\mAccess{M}ww'$.}}{}

\end{enumerate}
\end{defn}

As in other quantified logics, 
!!^{sentence}s are true at a world irrespective of assignment. 
Truth at a world is defined as in classical quantified modal logic. 

\begin{prop}
\Article{sentence} !!{sentence} is satisfied at a world relative to an
assignment if and only if it is satisfied at that world relative to to
any other assignment. That is, if $!A$ is !!a{sentence} then:
$$\qmSat{M}{!A}[w][s]\text{ iff }\qmSat{M}{!A}[w][s']\text{ for any }s,s'$$
\end{prop}

\begin{defn}[truth-at-w]
\Article{sentence} !!{sentence} $!A$ is \emph{true at a world} $w$
in model $\mModel M$ iff it is satisfied at $w$ relative to some assignment, 
or equivalently, iff it is satisfied at $w$ to every assignment.
We then simply write $\qmSat{M}{!A}[w]$. 
\end{defn}

\begin{ex}
The !!{sentence} $\lforall[x][\Box\lexists[y][\eq[x][y]]]$ 
(``Everything necessarily exists'') is false in some variable-domain 
model. For consider the model with:
\begin{itemize}
  \item $\mWorlds{M}=\{w,v\}$,
  \item $\mAccess{M}=\{\langle w,v\rangle\}$,
  \item $\mathcal{S}=\{1\}$,
  \item $\mDomain{M}(w)=\{1\}$ and $\mDomain{M}(v)=\varnothing$,
  \item $\mAssign{e}{M}$ any arbitrary interpretation of !!{constant},
  !!{predicate}, or !!{function} $e$.
\end{itemize}
Let $w$ $s$ be any !!{variable} assignment in $\mModel{M}$ such that
$s(x)=1$. Since $1$ is not in $\mDomain{M}(v)$ (in fact, nothing in is
$\mDomain{M}(v)$), there is no $\Subst{s}{m}{y}$ with $m$ in
$\mDomain{M}(v)$ such that $\Subst{s}{m}{y}=\Subst{s}{m}{x}=s(x)=1$.
Hence $\qmSat/{M}{x=y}[v][\Subst{s}{m}{x}]$ for any $m$ in
$\mDomain{M}(v)$. Therefore (by \olref{defn:satisfaction})
$\qmSat/{M}{\lexists[y][x=y]}[v][s]$. Since $v$ is accessible from
$w$, (by \olref{defn:satisfaction})
$\qmSat/{M}{\Box\lexists[y][x=y]}[w][s]$. Since $s(x)=1$
and $1\in \mDomain{M}(w)$, we have  
$\qmSat/{M}{\Box\lexists[y][x=y]}[w][\Subst{s}{x}{1}]$
with $1\in \mDomain{M}(w)$, therefore (by \olref{defn:satisfaction})
$\qmSat/{M}{\forall x\Box\lexists[y][x=y]}[w][\Subst{s}{x}{1}]$

All that the model requires is that one world (here, $w$) has
a domain that includes some object (here, $1$) and has access
to a world (here, $v$) whose domain does not include that object.
\end{ex}

\begin{rem}
Fixed-domain models are (in effect) a special case of variable-domain
models. Call a variable-domain \emph{fixed} iff its super-domain is
the domain of each world: $\mDomain{M}(w)=\mathcal{S}$ for every
$w\in\mWorlds{M}$. It's easy to see that $\mModel{M}$ corresponds to a
fixed-domain model $\mModel{M}^*$ with $\mDomain{M}^*=\mathcal{S}$, in
the sense that the same formulas are satisfied relative to the same
worlds and assignments in $\mModel{M}$ and in $\mModel{M}^*$.
\end{rem}

\end{document}
